% ============================================================
\chapter{Conteneurisation --- Docker pour le ML}
\label{ch:docker}
\index{Docker}
% ============================================================

\section{Pourquoi les conteneurs~?}
\label{sec:why-containers}

Le probl\`eme <<~\c{C}a marche sur ma machine~!~>> persiste m\^eme avec des
environnements virtuels. Les causes profondes~:
\begin{itemize}
  \item Diff\'erences de syst\`eme d'exploitation (Ubuntu vs Debian vs Alpine)
  \item Versions de biblioth\`eques syst\`eme (glibc, OpenSSL)
  \item Drivers GPU et versions de CUDA non align\'es
  \item Variables d'environnement, locales, fuseaux horaires
  \item Services externes (bases de donn\'ees, files de messages)
\end{itemize}

Un environnement virtuel isole les packages Python, mais \textbf{pas} le
syst\`eme d'exploitation ni les d\'ependances syst\`eme. Les conteneurs
r\'esolvent ce probl\`eme en encapsulant \emph{tout} l'environnement
d'ex\'ecution.

\begin{definition}[Conteneur]
Un \textbf{conteneur}\index{conteneur} est une unit\'e logicielle l\'eg\`ere
et portable qui empaquette le code, les d\'ependances et la configuration
dans un environnement isol\'e, partageant le noyau du syst\`eme h\^ote.
\end{definition}

\begin{remarque}
Contrairement \`a une machine virtuelle, un conteneur ne virtualise pas le
mat\'eriel~: il partage le noyau Linux de l'h\^ote. Cela le rend beaucoup
plus l\'eger (d\'emarrage en secondes, empreinte m\'emoire r\'eduite).
\end{remarque}

\subsection{Conteneurs vs Machines virtuelles}

\begin{center}
\begin{tikzpicture}[
  layer/.style={rectangle, draw, minimum width=5cm, minimum height=0.7cm,
    align=center, font=\scriptsize},
  title/.style={font=\small\bfseries, above=0.1cm}
]
  % VM side
  \node[layer, fill=gray!20] (hw1) {Infrastructure (mat\'eriel)};
  \node[layer, fill=orange!15, above=0cm of hw1] (hyp) {Hyperviseur};
  \node[layer, fill=blue!10, above=0cm of hyp, minimum width=1.5cm,
    xshift=-1.5cm] (os1) {OS invit\'e};
  \node[layer, fill=blue!10, above=0cm of hyp, minimum width=1.5cm,
    xshift=0cm] (os2) {OS invit\'e};
  \node[layer, fill=blue!10, above=0cm of hyp, minimum width=1.5cm,
    xshift=1.5cm] (os3) {OS invit\'e};
  \node[layer, fill=green!15, above=0cm of os1, minimum width=1.5cm,
    xshift=0cm] (app1) {App A};
  \node[layer, fill=green!15, above=0cm of os2, minimum width=1.5cm,
    xshift=0cm] (app2) {App B};
  \node[layer, fill=green!15, above=0cm of os3, minimum width=1.5cm,
    xshift=0cm] (app3) {App C};
  \node[title] at ([yshift=0.3cm]app2.north) {Machine Virtuelle};

  % Container side
  \node[layer, fill=gray!20, right=3cm of hw1] (hw2) {Infrastructure (mat\'eriel)};
  \node[layer, fill=orange!15, above=0cm of hw2] (hos) {OS h\^ote + Docker Engine};
  \node[layer, fill=green!15, above=0cm of hos, minimum width=1.5cm,
    xshift=-1.5cm] (c1) {App A};
  \node[layer, fill=green!15, above=0cm of hos, minimum width=1.5cm,
    xshift=0cm] (c2) {App B};
  \node[layer, fill=green!15, above=0cm of hos, minimum width=1.5cm,
    xshift=1.5cm] (c3) {App C};
  \node[title] at ([yshift=0.3cm]c2.north) {Conteneurs};
\end{tikzpicture}
\end{center}

\begin{center}
\begin{tabular}{@{}lcc@{}}
\toprule
\textbf{Crit\`ere} & \textbf{VM} & \textbf{Conteneur} \\
\midrule
D\'emarrage & Minutes & Secondes \\
Taille & Go & Mo \\
Isolation & Compl\`ete & Niveau processus \\
Performance & Overhead noyau & Quasi-native \\
Portabilit\'e & Image lourde & Image l\'eg\`ere \\
GPU passthrough & Complexe & NVIDIA Container Toolkit \\
\bottomrule
\end{tabular}
\end{center}

\section{Concepts Docker}
\label{sec:docker-concepts}
\index{Docker!concepts}

\begin{formules}[title=\textbf{Vocabulaire Docker essentiel}]
\begin{description}
  \item[Image]\index{Docker!image} Mod\`ele en lecture seule contenant le
    syst\`eme de fichiers, les d\'ependances et le code. Analogue \`a une
    <<~classe~>> en POO.
  \item[Conteneur]\index{Docker!conteneur} Instance en cours d'ex\'ecution
    d'une image. Analogue \`a un <<~objet~>> instanci\'e depuis la classe.
  \item[Dockerfile]\index{Dockerfile} Fichier texte d\'ecrivant les
    instructions pour construire une image, couche par couche.
  \item[Layer (couche)]\index{Docker!layer} Chaque instruction du Dockerfile
    cr\'ee une couche. Les couches sont cach\'ees et r\'eutilis\'ees.
  \item[Registry]\index{Docker!registry} D\'ep\^ot d'images (Docker Hub,
    GitHub Container Registry, AWS ECR).
  \item[Volume]\index{Docker!volume} M\'ecanisme de persistance des donn\'ees
    en dehors du conteneur.
  \item[Tag] \'Etiquette de version d'une image (ex.\ \cli{python:3.11-slim}).
\end{description}
\end{formules}

\section{Dockerfile pour projets ML}
\label{sec:dockerfile-ml}
\index{Dockerfile!ML}

\begin{codeblock}[title=\textbf{Dockerfile ML complet}]
\begin{minted}{dockerfile}
# Image de base avec Python
FROM python:3.11-slim AS base

# Metadonnees
LABEL maintainer="equipe-ml@example.com"
LABEL description="ML training and inference image"

# Variables d'environnement
ENV PYTHONDONTWRITEBYTECODE=1 \
    PYTHONUNBUFFERED=1 \
    PIP_NO_CACHE_DIR=1

# Dependances systeme
RUN apt-get update && apt-get install -y --no-install-recommends \
    build-essential \
    curl \
    && rm -rf /var/lib/apt/lists/*

# Creer un utilisateur non-root
RUN useradd --create-home mluser
WORKDIR /home/mluser/app

# Installer les dependances Python d'abord (cache Docker)
COPY requirements.txt .
RUN pip install --no-cache-dir -r requirements.txt

# Copier le code source
COPY src/ ./src/
COPY configs/ ./configs/

# Changer de proprietaire et d'utilisateur
RUN chown -R mluser:mluser /home/mluser/app
USER mluser

# Port pour l'API (si applicable)
EXPOSE 8000

# Commande par defaut
CMD ["python", "src/train.py", "--config", "configs/config.yaml"]
\end{minted}
\end{codeblock}

\begin{bonnepratique}[title=\textbf{Ordre des instructions et cache}]
Docker met en cache chaque couche. Si une couche change, toutes les couches
suivantes sont reconstruites. R\`egle d'or~:
\begin{enumerate}
  \item Placer les \'el\'ements qui changent \textbf{rarement} en haut
    (OS, d\'ependances syst\`eme).
  \item Copier \file{requirements.txt} \textbf{avant} le code source.
  \item Copier le code source en \textbf{dernier}.
\end{enumerate}
Ainsi, modifier le code ne d\'eclenche pas la r\'einstallation des packages.
\end{bonnepratique}

\subsection{Dockerfile avec support CUDA}

\begin{codeblock}[title=\textbf{Dockerfile GPU avec CUDA}]
\begin{minted}{dockerfile}
# Image de base NVIDIA CUDA
FROM nvidia/cuda:12.1.0-cudnn8-runtime-ubuntu22.04

ENV DEBIAN_FRONTEND=noninteractive \
    PYTHONDONTWRITEBYTECODE=1 \
    PYTHONUNBUFFERED=1

# Installer Python
RUN apt-get update && apt-get install -y --no-install-recommends \
    python3.11 python3-pip python3.11-venv \
    && ln -s /usr/bin/python3.11 /usr/bin/python \
    && rm -rf /var/lib/apt/lists/*

WORKDIR /app

COPY requirements.txt .
RUN pip install --no-cache-dir -r requirements.txt

COPY src/ ./src/
COPY configs/ ./configs/

RUN useradd --create-home mluser && chown -R mluser:mluser /app
USER mluser

CMD ["python", "src/train.py"]
\end{minted}
\end{codeblock}

\section{Commandes Docker essentielles}
\label{sec:docker-commands}

\begin{codeblock}[title=\textbf{Cycle de vie d'un conteneur}]
\begin{minted}{bash}
# Construire une image
docker build -t mon-projet-ml:v1.0 .

# Lancer un conteneur
docker run --name ml-train mon-projet-ml:v1.0

# Lancer en mode interactif
docker run -it mon-projet-ml:v1.0 bash

# Lancer avec un volume (donnees persistantes)
docker run -v $(pwd)/data:/app/data mon-projet-ml:v1.0

# Lancer avec GPU
docker run --gpus all mon-projet-ml:v1.0

# Lister les conteneurs
docker ps -a

# Voir les logs
docker logs ml-train

# Arreter et supprimer
docker stop ml-train && docker rm ml-train

# Lister et supprimer les images
docker images
docker rmi mon-projet-ml:v1.0
\end{minted}
\end{codeblock}

\section{Multi-stage builds}
\label{sec:multi-stage}
\index{Docker!multi-stage build}

Les \textbf{multi-stage builds} permettent de s\'eparer l'environnement de
construction de l'environnement d'ex\'ecution, r\'eduisant drastiquement la
taille de l'image finale.

\begin{codeblock}[title=\textbf{Multi-stage build pour ML}]
\begin{minted}{dockerfile}
# === Stage 1 : construction ===
FROM python:3.11-slim AS builder

RUN apt-get update && apt-get install -y --no-install-recommends \
    build-essential gfortran libopenblas-dev \
    && rm -rf /var/lib/apt/lists/*

WORKDIR /build
COPY requirements.txt .
RUN pip install --no-cache-dir --prefix=/install -r requirements.txt

# === Stage 2 : execution ===
FROM python:3.11-slim AS runtime

# Copier seulement les packages installes
COPY --from=builder /install /usr/local

RUN useradd --create-home mluser
WORKDIR /home/mluser/app

COPY src/ ./src/
COPY configs/ ./configs/
COPY models/ ./models/

RUN chown -R mluser:mluser /home/mluser/app
USER mluser

EXPOSE 8000
CMD ["python", "src/serve.py"]
\end{minted}
\end{codeblock}

\begin{center}
\begin{tabular}{@{}lcc@{}}
\toprule
\textbf{Approche} & \textbf{Taille image} & \textbf{S\'ecurit\'e} \\
\midrule
Single-stage (python:3.11) & $\sim$1.2~Go & Compilateurs pr\'esents \\
Single-stage (python:3.11-slim) & $\sim$800~Mo & Am\'elior\'ee \\
Multi-stage (slim $\to$ slim) & $\sim$400~Mo & Pas de compilateurs \\
\bottomrule
\end{tabular}
\end{center}

\section{Docker Compose pour environnements multi-services}
\label{sec:docker-compose}
\index{Docker Compose}

En ML, on a souvent besoin de plusieurs services~: l'application, une base
de donn\'ees pour les m\'etadonn\'ees, un serveur MLflow pour le suivi des
exp\'eriences, etc. \textbf{Docker Compose}\index{Docker Compose} orchestre
ces services dans un seul fichier.

\begin{codeblock}[title=\textbf{docker-compose.yml --- Stack ML compl\`ete}]
\begin{minted}{yaml}
version: "3.9"

services:
  # --- Application ML ---
  ml-app:
    build:
      context: .
      dockerfile: Dockerfile
    ports:
      - "8000:8000"
    volumes:
      - ./data:/app/data
      - ./models:/app/models
    environment:
      - MLFLOW_TRACKING_URI=http://mlflow:5000
      - DATABASE_URL=postgresql://user:pass@db:5432/mldb
    depends_on:
      - db
      - mlflow
    deploy:
      resources:
        reservations:
          devices:
            - capabilities: [gpu]

  # --- Base de donnees PostgreSQL ---
  db:
    image: postgres:16-alpine
    environment:
      POSTGRES_USER: user
      POSTGRES_PASSWORD: pass
      POSTGRES_DB: mldb
    volumes:
      - pgdata:/var/lib/postgresql/data
    ports:
      - "5432:5432"

  # --- MLflow Tracking Server ---
  mlflow:
    image: ghcr.io/mlflow/mlflow:v2.10.0
    command: >
      mlflow server
      --backend-store-uri postgresql://user:pass@db:5432/mldb
      --default-artifact-root /mlflow/artifacts
      --host 0.0.0.0
      --port 5000
    ports:
      - "5000:5000"
    volumes:
      - mlflow-artifacts:/mlflow/artifacts
    depends_on:
      - db

volumes:
  pgdata:
  mlflow-artifacts:
\end{minted}
\end{codeblock}

\begin{codeblock}[title=\textbf{Commandes Docker Compose}]
\begin{minted}{bash}
# Demarrer tous les services
docker compose up -d

# Voir les logs de tous les services
docker compose logs -f

# Voir les logs d'un service specifique
docker compose logs -f ml-app

# Arreter tous les services
docker compose down

# Arreter et supprimer les volumes
docker compose down -v

# Reconstruire les images
docker compose build --no-cache

# Lancer une commande dans un service
docker compose exec ml-app python src/train.py
\end{minted}
\end{codeblock}

\section{Le fichier .dockerignore}
\label{sec:dockerignore}
\index{.dockerignore}

Le fichier \file{.dockerignore} emp\^eche l'envoi de fichiers inutiles
au d\'emon Docker lors du build. Sans lui, Docker copie \emph{tout} le
contexte, y compris les donn\'ees, les mod\`eles et les environnements
virtuels.

\begin{codeblock}[title=\textbf{.dockerignore pour projet ML}]
\begin{minted}{text}
# Environnements virtuels
.venv/
venv/
__pycache__/
*.pyc

# Donnees volumineuses
data/raw/
data/processed/
*.csv
*.parquet

# Modeles sauvegardes
models/*.joblib
models/*.pkl
*.onnx

# Git et IDE
.git/
.gitignore
.vscode/
.idea/

# Notebooks
*.ipynb
.ipynb_checkpoints/

# Docker
Dockerfile
docker-compose.yml
.dockerignore

# Documentation
*.md
docs/
\end{minted}
\end{codeblock}

\section{Bonnes pratiques Docker pour le ML}
\label{sec:docker-best-practices}

\begin{bonnepratique}[title=\textbf{R\`egles d'or Docker en ML}]
\begin{enumerate}
  \item \textbf{Images l\'eg\`eres}~: utiliser \cli{python:3.11-slim} ou
    \cli{alpine}, pas \cli{python:3.11} (\'economie de 600~Mo).
  \item \textbf{Utilisateur non-root}~: toujours cr\'eer et utiliser un
    utilisateur d\'edi\'e (\cli{USER mluser}).
  \item \textbf{Exploiter le cache}~: copier \file{requirements.txt} avant
    le code source.
  \item \textbf{Multi-stage builds}~: s\'eparer construction et ex\'ecution.
  \item \textbf{.dockerignore}~: exclure donn\'ees, mod\`eles, \file{.git/}.
  \item \textbf{Tags explicites}~: \cli{python:3.11-slim}, jamais
    \cli{python:latest}.
  \item \textbf{Un processus par conteneur}~: ne pas lancer la base de
    donn\'ees et l'application dans le m\^eme conteneur.
  \item \textbf{Healthchecks}~: ajouter \cli{HEALTHCHECK} pour les services
    de serving.
\end{enumerate}
\end{bonnepratique}

\begin{antipattern}[title=\textbf{Anti-patterns Docker en ML}]
\begin{itemize}
  \item \textbf{Images g\'eantes}~: inclure les compilateurs, les donn\'ees
    d'entra\^inement et les checkpoints dans l'image (r\'esultat~: 10+~Go).
  \item \textbf{Ex\'ecution en root}~: lancer le conteneur en tant que root
    expose l'h\^ote en cas de vuln\'erabilit\'e.
  \item \textbf{Secrets dans l'image}~: mettre des cl\'es API, mots de passe
    ou tokens dans le Dockerfile ou dans les couches de l'image. Utiliser
    plut\^ot des variables d'environnement ou Docker Secrets.
  \item \textbf{Tag \cli{latest}}~: ne permet pas de reproduire un build.
    Toujours \'epingler une version pr\'ecise.
  \item \textbf{Ignorer le .dockerignore}~: envoyer des Go de donn\'ees au
    d\'emon Docker \`a chaque build.
  \item \textbf{Un seul gros conteneur}~: mettre l'app, la BDD et MLflow
    dans le m\^eme conteneur.
  \item \textbf{Pas de healthcheck}~: impossible de savoir si le service
    est r\'eellement fonctionnel.
\end{itemize}
\end{antipattern}

\section{Registries et publication d'images}
\label{sec:registries}
\index{Docker!registry}

\begin{codeblock}[title=\textbf{Publier une image sur Docker Hub}]
\begin{minted}{bash}
# Se connecter
docker login

# Taguer l'image
docker tag mon-projet-ml:v1.0 monuser/mon-projet-ml:v1.0

# Pousser l'image
docker push monuser/mon-projet-ml:v1.0

# Tirer l'image (sur une autre machine)
docker pull monuser/mon-projet-ml:v1.0
\end{minted}
\end{codeblock}

\begin{codeblock}[title=\textbf{Utiliser GitHub Container Registry}]
\begin{minted}{bash}
# Se connecter avec un token GitHub
echo $GITHUB_TOKEN | docker login ghcr.io -u USERNAME --password-stdin

# Taguer et pousser
docker tag mon-projet-ml:v1.0 ghcr.io/USERNAME/mon-projet-ml:v1.0
docker push ghcr.io/USERNAME/mon-projet-ml:v1.0
\end{minted}
\end{codeblock}

\section{Mini-projet~: conteneuriser le projet ML du cours}
\label{sec:mini-project-ch7}

\begin{miniprojet}[title=\textbf{Mini-projet 7~: Conteneurisation compl\`ete}]
En reprenant le projet ML structur\'e des chapitres pr\'ec\'edents~:
\begin{enumerate}
  \item \'Ecrivez un \file{Dockerfile} multi-stage pour le projet~:
    \begin{itemize}
      \item Stage 1~: installer les d\'ependances et compiler
      \item Stage 2~: copier uniquement les packages et le code
    \end{itemize}
  \item Cr\'eez un \file{.dockerignore} adapt\'e.
  \item Construisez l'image et v\'erifiez sa taille
    (\cli{docker images}).
  \item Lancez un entra\^inement dans le conteneur avec un volume
    pour les donn\'ees~: \cli{docker run -v \$(pwd)/data:/app/data ...}
  \item \'Ecrivez un \file{docker-compose.yml} avec trois services~:
    \begin{itemize}
      \item \cli{ml-app}~: votre application ML
      \item \cli{db}~: PostgreSQL pour MLflow
      \item \cli{mlflow}~: serveur de suivi des exp\'eriences
    \end{itemize}
  \item V\'erifiez que \cli{docker compose up} d\'emarre tout
    correctement et que MLflow est accessible sur
    \cli{http://localhost:5000}.
\end{enumerate}
\end{miniprojet}

\section{Exercices}
\label{sec:exercises-ch7}

\begin{exercice}[$\bigstar$ --- Premier Dockerfile]
\'Ecrivez un \file{Dockerfile} pour un script Python simple qui charge un
dataset CSV avec \pkg{pandas} et affiche des statistiques descriptives.
Construisez l'image, lancez le conteneur et v\'erifiez la sortie.
Comparez la taille de l'image avec \cli{python:3.11} vs
\cli{python:3.11-slim}.
\end{exercice}

\begin{exercice}[$\bigstar\bigstar$ --- Optimisation d'image]
Partez d'un Dockerfile na\"if (image \cli{python:3.11}, pas de
\file{.dockerignore}, copie de tout le r\'epertoire, ex\'ecution en root).
Am\'eliorez-le it\'erativement~:
\begin{enumerate}
  \item Passez \`a \cli{python:3.11-slim} --- notez la r\'eduction de taille
  \item Ajoutez un \file{.dockerignore} --- mesurez le temps de build
  \item R\'eordonnez les instructions pour exploiter le cache
  \item Passez en multi-stage build
  \item Ajoutez un utilisateur non-root
\end{enumerate}
Documentez la taille de l'image et le temps de build \`a chaque \'etape.
\end{exercice}

\begin{exercice}[$\bigstar\bigstar\bigstar$ --- Stack ML compl\`ete avec Docker Compose]
Cr\'eez un \file{docker-compose.yml} complet pour un projet ML incluant~:
\begin{enumerate}
  \item Un service d'entra\^inement (GPU si disponible)
  \item Un service d'inf\'erence FastAPI (port 8000)
  \item Un serveur MLflow (port 5000)
  \item Une base de donn\'ees PostgreSQL
  \item Un dashboard Streamlit (port 8501)
  \item Des healthchecks pour chaque service
  \item Des volumes pour la persistance des donn\'ees et des mod\`eles
\end{enumerate}
V\'erifiez que tous les services communiquent correctement~: l'entra\^inement
loggue dans MLflow, l'API charge le mod\`ele depuis le volume partag\'e, et
le dashboard interroge l'API.
\end{exercice}

\section{Aide-m\'emoire}

\begin{cheatsheet}[title=\textbf{R\'esum\'e du chapitre 7}]
\begin{tabular}{@{}p{4cm}p{8.5cm}@{}}
\toprule
\textbf{Concept} & \textbf{Commande / Point cl\'e} \\
\midrule
Construire & \cli{docker build -t name:tag .} \\
Lancer & \cli{docker run -v data:/app/data name:tag} \\
GPU & \cli{docker run --gpus all name:tag} \\
Multi-stage & \cli{FROM ... AS builder} puis \cli{COPY --from=builder} \\
Compose up & \cli{docker compose up -d} \\
Compose down & \cli{docker compose down -v} \\
Cache optimal & \file{requirements.txt} avant code source \\
S\'ecurit\'e & \cli{USER mluser}, jamais de secrets dans l'image \\
Image l\'eg\`ere & \cli{python:3.11-slim} + multi-stage \\
\bottomrule
\end{tabular}
\end{cheatsheet}
